\chapter*{Abstract in lingua italiana}


L'identificazione dei tumori nei tessuti si basa attualmente sull'esame istopatologico, durante il quale i campioni di biopsia vengono preparati, trattati con coloranti istologici e valutati da patologi esperti. Sebbene questo approccio rappresenti il riferimento clinico valido, richiede una preparazione specifica del campione, competenze specialistiche e tempi di analisi prolungati. Inoltre, la valutazione può includere un'influenza soggettiva/specialistica, che assume un ruolo soprattutto quando i campioni sono particolarmente eterogenei, infiammati o presentano stroma reattivo. Lo sviluppo di metodi quantitativi e riproducibili, in grado di supportare e, in futuro, di automatizzare parzialmente l'ispezione istologica, rappresenta una direzione importante per la ricerca biomedica.

 La spettroscopia Raman emerge come una tecnica promettente in questo scenario, poiché fornisce informazioni chimicamente specifiche, non richiede l'uso di marcatori e non è distruttiva. L'alto contenuto informativo degli spettri Raman, dovuto alla complessità dei campioni biologici, rende necessario ricorrere a strumenti computazionali e di apprendimento automatico adatti all'interpretazione e alla classificazione di dati complessi e ad alta dimensionalità. 
 
 Questa tesi sviluppa una pipeline di apprendimento automatico per l'analisi di mappe iperspettrali di Raman spontaneo acquisite da biopsie di carcinoma squamoso della testa e del collo (HNSCC) derivate dal paziente. Il lavoro combina l'Analisi delle Componenti Principali (PCA) e l'allenamento di una Rete Neurale Artificiale (ANN), rispettivamente, per la riduzione della dimensionalità e la classificazione supervisionata dei dati Raman, finalizzate a identificare il tessuto tumorale. Nel complesso, i risultati evidenziano il potenziale della spettroscopia Raman nel cogliere le differenze biochimiche e mostrano come l'apprendimento automatico possa supportare lo sviluppo di procedure computazionali per la caratterizzazione dei tessuti tumorali

\vspace{1em}
\noindent\textbf{Parole chiave}: spettroscopia Raman, mappe iperspettrali, apprendimento automatico, caratterizzazione dei tessuti tumorali.